\chapter{Chargement des donn\'ees}
\label{ch:data-loading}

\begin{quote}
\emph{<< On ne peut pas analyser ce qu'on ne peut pas charger. >>}
\end{quote}

% --- hook ---
Un CSV n'est pas un format. C'est une convention floue, et chaque équipe qui en a écrit a légèrement modifié la convention. Les Allemands utilisent la virgule comme séparateur décimal et le point-virgule comme séparateur de colonnes. Les Américains font l'inverse. Les fichiers Excel exportés en CSV portent souvent une signature BOM invisible qui casse tout. Les anciens systèmes bancaires écrivent encore en latin-1 et lèvent une exception silencieuse à la première signature wolof avec un accent. Ce chapitre couvre les six sources de données les plus fréquentes en pratique (CSV, Excel, JSON, SQL, API, web scraping), avec leurs pièges connus --- ceux qui passent en revue de code, et ceux qu'on ne voit qu'en production.
\section{Charger des fichiers CSV}

CSV (comma-separated values) est le format d'\'echange de donn\'ees le plus courant. Pandas le g\`ere avec \texttt{read\_csv}, qui a plus de 50 param\`etres pour traiter les particularit\'es du monde r\'eel.

\begin{minted}{python}
import pandas as pd

# Chargement basique
df = pd.read_csv("titanic.csv")
print(df.shape)
print(df.head())
\end{minted}

\subsection{Probl\`emes CSV courants}

\begin{minted}{python}
# Probl\`eme 1 : mauvais d\'elimiteur
# Certains CSV europ\'eens utilisent des points-virgules
df_euro = pd.read_csv("data_eu.csv", sep=";", decimal=",")

# Probl\`eme 2 : encodages
df_fr = pd.read_csv("french_data.csv", encoding="latin-1")

# Probl\`eme 3 : ent\^etes d\'esordonn\'ees
df = pd.read_csv("messy.csv", header=0, skiprows=[1, 2])
df.columns = df.columns.str.strip().str.lower().str.replace(" ", "_")
print(df.columns.tolist())
\end{minted}

\begin{preprocesstip}
Standardisez toujours les noms de colonnes imm\'ediatement apr\`es chargement : minuscules, underscores, pas d'espaces. \c{C}a \'evite des bugs subtils dus \`a une capitalisation incoh\'erente.
\end{preprocesstip}

\subsection{Optimiser le chargement CSV}

\begin{minted}{python}
# Charger uniquement les colonnes n\'ecessaires (\'economie m\'emoire)
cols = ["PassengerId", "Survived", "Pclass", "Sex", "Age", "Fare"]
df = pd.read_csv("titanic.csv", usecols=cols)

# Sp\'ecifier les dtypes pour \'economiser la m\'emoire
dtypes = {"Survived": "int8", "Pclass": "int8", "Sex": "category"}
df = pd.read_csv("titanic.csv", usecols=cols, dtype=dtypes)
print(df.memory_usage(deep=True))

# Charger par morceaux pour les gros fichiers
chunks = pd.read_csv("large_file.csv", chunksize=50_000)
result = pd.concat([chunk[chunk["Age"] > 30] for chunk in chunks])
\end{minted}

\begin{warningbox}
Ne chargez jamais un CSV de plusieurs gigaoctets enti\`erement en m\'emoire d'un coup. Utilisez \texttt{chunksize}, ou passez \`a Parquet/Feather. Un CSV de 2\,Go fait g\'en\'eralement 400\,Mo en Parquet avec des temps de chargement bien plus courts.
\end{warningbox}

% ============================================================
\section{Charger des fichiers Excel}

\begin{minted}{python}
# Chargement Excel basique (n\'ecessite openpyxl)
df = pd.read_excel("survey_results.xlsx", sheet_name="Sheet1")

# Charger toutes les feuilles dans un dictionnaire
all_sheets = pd.read_excel("multi_sheet.xlsx", sheet_name=None)
for name, sheet_df in all_sheets.items():
    print(f"Sheet '{name}': {sheet_df.shape}")

# Sauter ent\^etes et pieds de page
df = pd.read_excel("report.xlsx", skiprows=3, skipfooter=2,
                    engine="openpyxl")
\end{minted}

\begin{datatip}
Les fichiers Excel pr\'eservent le formatage mais perdent la pr\'ecision de type. Les dates peuvent \^etre charg\'ees comme cha\^ines, les pourcentages comme floats entre 0 et 1, les devises comme nombres bruts. V\'erifiez toujours les dtypes apr\`es chargement Excel.
\end{datatip}

% ============================================================
\section{Charger des donn\'ees JSON}

\begin{minted}{python}
import json

# JSON plat
df = pd.read_json("flat_data.json")

# JSON imbriqu\'e (courant depuis les API)
with open("nested_data.json") as f:
    raw = json.load(f)

# Utiliser json_normalize pour aplatir
from pandas import json_normalize

df = json_normalize(raw["results"],
                    record_path="measurements",
                    meta=["station_id", "station_name"])
print(df.head())
\end{minted}

\begin{minted}{python}
# Exemple r\'eel : charger un JSON de type Gapminder
import requests

url = "https://api.worldbank.org/v2/country/BEN/indicator/SP.POP.TOTL?format=json&per_page=60"
response = requests.get(url)
data = response.json()

# Les vraies donn\'ees sont dans data[1]
df_wb = pd.DataFrame(data[1])
print(df_wb[["date", "value"]].head(10))
\end{minted}

% ============================================================
\section{Charger depuis des bases SQL}

\begin{minted}{python}
from sqlalchemy import create_engine

# SQLite (fichier, pas de serveur)
engine = create_engine("sqlite:///local_data.db")

# Charger toute la table
df = pd.read_sql("SELECT * FROM patients", engine)

# Charger avec filtrage (pousser le calcul en base)
query = """
SELECT patient_id, age, diagnosis, lab_value
FROM patients
WHERE age > 40
  AND diagnosis IS NOT NULL
ORDER BY lab_value DESC
LIMIT 10000
"""
df = pd.read_sql(query, engine)
print(df.shape)
\end{minted}

\begin{preprocesstip}
En chargeant depuis SQL, poussez autant de filtrage que possible dans la requ\^ete. Charger 10 millions de lignes dans pandas puis filtrer est bien plus lent que laisser la base filtrer d'abord.
\end{preprocesstip}

% ============================================================
\section{Charger depuis des API}

\begin{minted}{python}
import requests
import pandas as pd

# Exemple d'API REST : Open Meteo (m\'et\'eo)
url = "https://archive-api.open-meteo.com/v1/archive"
params = {
    "latitude": 6.36,    # Cotonou
    "longitude": 2.43,
    "start_date": "2024-01-01",
    "end_date": "2024-12-31",
    "daily": "temperature_2m_mean,precipitation_sum",
}

response = requests.get(url, params=params)
data = response.json()

df_weather = pd.DataFrame({
    "date": data["daily"]["time"],
    "temp_mean": data["daily"]["temperature_2m_mean"],
    "precip_mm": data["daily"]["precipitation_sum"],
})
df_weather["date"] = pd.to_datetime(df_weather["date"])
print(df_weather.head())
\end{minted}

\begin{minted}{python}
# Pagination : beaucoup d'API renvoient les donn\'ees par pages
all_records = []
page = 1

while True:
    resp = requests.get(f"{base_url}?page={page}&per_page=100")
    records = resp.json()["results"]
    if not records:
        break
    all_records.extend(records)
    page += 1

df = pd.DataFrame(all_records)
print(f"Total records: {len(df)}")
\end{minted}

% ============================================================
\section{Bases du web scraping}

\begin{minted}{python}
import requests
from bs4 import BeautifulSoup
import pandas as pd

# Scraper un tableau HTML
url = "https://en.wikipedia.org/wiki/List_of_countries_by_GDP_(nominal)"
response = requests.get(url)
soup = BeautifulSoup(response.content, "html.parser")

# pandas peut lire les tableaux HTML directement
tables = pd.read_html(url)
print(f"Found {len(tables)} tables on the page")
gdp_df = tables[2]  # choisir le bon tableau par index
print(gdp_df.head())
\end{minted}

\begin{warningbox}
Le web scraping peut violer les conditions d'utilisation d'un site. V\'erifiez toujours \texttt{robots.txt} et les ToS avant de scraper. Pour la recherche, pr\'ef\'erez les API officielles et les portails de donn\'ees ouvertes.
\end{warningbox}

% ============================================================
\section{Comparer les formats : performance}

\begin{minted}{python}
import time

# Comparer les temps de chargement CSV vs Parquet
# D'abord, cr\'eer une version Parquet
df = pd.read_csv("large_dataset.csv")
df.to_parquet("large_dataset.parquet")

# Chronom\'etrer le chargement CSV
start = time.time()
df_csv = pd.read_csv("large_dataset.csv")
csv_time = time.time() - start

# Chronom\'etrer le chargement Parquet
start = time.time()
df_pq = pd.read_parquet("large_dataset.parquet")
pq_time = time.time() - start

print(f"CSV load:     {csv_time:.2f}s")
print(f"Parquet load: {pq_time:.2f}s")
print(f"Speedup:      {csv_time / pq_time:.1f}x")
\end{minted}

% ============================================================
\section{Exercices}

\begin{exercise}
\begin{enumerate}
  \item Chargez le CSV Titanic avec des dtypes optimis\'es. Comparez l'usage m\'emoire avant et apr\`es optimisation.
  \item Chargez des donn\'ees depuis l'API Banque Mondiale pour trois pays africains (B\'enin, Nigeria, Ghana) sur l'indicateur \texttt{SP.POP.TOTL} (population totale). Combinez en un seul DataFrame.
  \item Chargez un fichier Excel \`a plusieurs feuilles. \'Ecrivez une fonction qui renvoie un dictionnaire associant le nom de feuille \`a sa shape.
  \item Scrapez un tableau HTML depuis Wikipedia et nettoyez les noms de colonnes (minuscules, underscores, pas de caract\`eres sp\'eciaux).
  \item Convertissez le dataset Adult Census UCI de CSV vers Parquet. Comparez tailles de fichiers et temps de chargement.
\end{enumerate}
\end{exercise}

% ============================================================
\section{R\'esum\'e du chapitre}

\begin{itemize}
  \item Pandas fournit des lecteurs pour CSV, Excel, JSON, SQL, Parquet, HTML, et plus.
  \item Le chargement CSV demande attention aux d\'elimiteurs, encodages, ent\^etes et dtypes.
  \item Le JSON issu d'API a souvent besoin de \texttt{json\_normalize} pour aplatir les structures imbriqu\'ees.
  \item Les chargements SQL devraient pousser le filtrage dans la requ\^ete, pas dans pandas.
  \item Les API demandent de g\'erer pagination, rate limits, et authentification.
  \item Parquet est 3 \`a 5 fois plus rapide que CSV au chargement et utilise 5 fois moins d'espace disque.
\end{itemize}
